% ==========================================================================
%  Chapter 102 -- Cultural-Algorithm Metaheuristics Module (Extension Plan, §M)
% ==========================================================================

\chapter{A Cultural-Algorithm Metaheuristics Module}
\label{ch:cultural-algorithm-module}

\paragraph{Normative status.}
Phase~M of the branch-selection extension
(Cap.~\ref{ch:branch-selection-reversals}), and the reusable substrate for
the Phase-P3 solver meta-tuner. Implemented header-only under
\path{src/algorithms/} and verified by ctest
\texttt{algorithms\_cultural\_algorithm}
(\path{tests/test_algorithms_cultural_algorithm.cpp}). Status:
\texttt{closed\_with\_runtime\_evidence} (Gate~G3): the module compiles in
the dependency-free \texttt{NONE} tier and is bit-for-bit deterministic
under a fixed seed.

\section{Why a Cultural Algorithm}

The solver exposes a handful of continuation controls (initial \(\mu\),
growth/drop factors, stagnation window, predictor scale, acceptance
floor, reversal subdivision). Their good values are coupled and
problem-dependent, and the objective --- ``converge the reversal without
a force spike'' --- is cheap to evaluate by replaying a short window but
expensive to differentiate. This is the natural habitat of a
population-based metaheuristic. A \emph{Cultural Algorithm} (CA)
\cite{Reynolds1994Cultural,ReynoldsPeng2004} is chosen over a plain
genetic algorithm because it adds a \emph{belief space}: knowledge
distilled from the best individuals is fed back to bias the search,
which accelerates convergence on the smooth, low-dimensional,
box-bounded landscapes that the solver's parameters inhabit. The same
machinery is the CA that tunes tuned-inerter dampers in
\cite{LaraValenciaEchavarriaValencia2024}; building it as a general
module rather than a solver-specific script is deliberate reuse.

\section{The dual-inheritance model}

A CA evolves two coupled populations \cite{Reynolds1994Cultural}:
%
\begin{itemize}
  \item the \emph{population space} --- the individuals, evolving by
        selection and variation (the ``genetic'' layer);
  \item the \emph{belief space} --- knowledge accumulated across
        generations (the ``cultural'' layer).
\end{itemize}
%
They communicate through two functions. The \emph{acceptance} function
selects the elite that updates the belief space each generation; the
\emph{influence} function lets the belief space bias the variation of the
next generation. Reynolds catalogues several \emph{knowledge sources};
this module implements the two that matter for smooth box-bounded search
and leaves the rest as documented extension points:
%
\begin{description}
  \item[Normative knowledge.] The per-gene interval
        \([\ell_j,h_j]\) spanned by the accepted elite. Influence
        resamples a gene uniformly inside its normative interval, pulling
        offspring toward the region the elite occupies.
  \item[Situational knowledge.] The best exemplar
        \(\mathbf{u}^\star\) seen so far. Influence nudges a gene toward
        \(\mathbf{u}^\star\) with a Gaussian step, exploiting the current
        champion.
  \item[Domain / History / Topographic.] Not needed here; the
        \texttt{KnowledgeSource} concept documents the two-method contract
        (\texttt{update(elite,space)}, \texttt{influence(genome,rng,space)})
        so they can be added without touching the belief space.
\end{description}

\section{Architecture (concepts and policies, no CRTP)}

The module is std-only and follows the library's compile-time-policy
idiom (Cap.~\ref{ch:steppable_solver_concept} style), namespace
\texttt{fall\_n::algorithms}:
%
\begin{itemize}
  \item \texttt{optimization/Concepts.hh}: \texttt{ObjectiveFunction}
        (over \texttt{std::span<const double>}, convention
        \emph{maximise}) and \texttt{SearchSpace}
        (\texttt{dimension/lower/upper/clamp/sample}).
  \item \texttt{optimization/BoundedSearchSpace.hh}: the box
        \([\boldsymbol\ell,\mathbf{h}]\); \texttt{static\_assert}s the
        \texttt{SearchSpace} concept.
  \item \texttt{cultural/\{Individual,KnowledgeSources,BeliefSpace,
        PopulationSpace,CulturalAlgorithm\}.hh}.
        \texttt{BeliefSpace<KS...>} holds the knowledge sources in a
        \texttt{std::tuple} and fans \texttt{accept}/\texttt{influence}
        over them with a fold expression, so adding a source is a
        template-argument change.
\end{itemize}
%
Randomness flows from a single \texttt{std::mt19937\_64} seeded by
\texttt{CulturalConfig::seed}; every draw --- initial sampling, tournament,
mutation, influence --- is taken from it, which is what makes a run
reproducible. \texttt{maximize} takes the RNG type as a defaulted method
template parameter (the class template ends in the knowledge-source pack,
so the RNG cannot follow it there), keeping injectability without
breaking the pack-last rule.

\section{The generational loop}

\begin{enumerate}
  \item Sample and evaluate the initial population from the search space.
  \item Rank by descending fitness.
  \item \(\texttt{belief.accept}\): the top-fraction elite updates the
        normative and situational sources.
  \item Build the next generation: carry the incumbent best (elitism),
        then fill by binary-tournament selection followed by variation
        --- with probability \(p_\text{infl}\) the belief space biases the
        child, otherwise a Gaussian mutation scaled by the box extent ---
        and clamp into the box.
  \item Track the best; stop on \texttt{max\_generations},
        \texttt{target\_fitness}, or a stagnation window.
\end{enumerate}
%
Elitism guarantees the best-fitness history is monotone
non-decreasing, which the test asserts.

\section{Verification (ctest \texttt{algorithms\_cultural\_algorithm})}

Standard benchmarks, each maximised as the negated cost, fixed seeds:
%
\begin{center}
\begin{tabular}{@{}lll@{}}
\toprule
Problem            & Best fitness reached & Meaning \\
\midrule
Sphere 5D          & \(-5.5\times10^{-3}\) & near the global optimum \(0\) \\
Rosenbrock 2D      & \(-3.5\times10^{-3}\) & inside the curved valley, near \((1,1)\) \\
Rastrigin 5D       & \(-3.0\)              & deep below the \(\sim\!-50\) random mean \\
\bottomrule
\end{tabular}
\end{center}
%
plus: the best genome respects the box on every problem, the
best-fitness history is non-decreasing (elitism), and two runs with the
same seed are \emph{bit-identical}. Thresholds are set below the achieved
values with margin so the deterministic test also guards against
regressions. Zero failures.

\section{The two applications}

\paragraph{Solver meta-tuning (Phase~P3).} The genome is the solver's
control vector; the objective replays a short window across the first
reversal from a captured checkpoint (Cap.~\ref{ch:deflated-continuation}
supplies \texttt{capture\_state}/\texttt{restore\_state}) and scores by
converged fraction minus a spike penalty. The search space is the box of
admissible controls; the CA returns a genome that is written back with
\texttt{set\_config}/\texttt{set\_regularization}. Material-latched genes
(closure stiffness fraction, anchor cap) are read once into
\texttt{static const} and so are tuned only by an outer per-process
sweep, never inside the replay loop --- a limitation recorded honestly.

\paragraph{Control-device design (future).} The identical
\texttt{maximize} drives tuned-mass/tuned-inerter damper design: the
genome carries device parameters, the objective installs them through
\texttt{DynamicAnalysis::set\_internal\_residual\_hook}, runs a ground
motion from a checkpoint, and returns \(-\,\text{peak drift}\). No FEM
dependency enters the module; the coupling lives entirely in the
caller-supplied objective. This mirrors, and will reuse, the CA of
\cite{LaraValenciaEchavarriaValencia2024}.
